\documentclass[11pt]{article}

\usepackage[a4paper,margin=25mm]{geometry}
\usepackage[T1]{fontenc}
\usepackage{lmodern}
\usepackage{amsmath,amssymb,amsthm}
\usepackage{enumitem}
\usepackage{hyperref}

\setlist[enumerate]{itemsep=0.35em,topsep=0.4em}
\setlist[itemize]{itemsep=0.25em,topsep=0.4em}
\newcommand{\Conf}{\mathsf{Conf}}
\newcommand{\Act}{\mathsf{Act}}
\newcommand{\eps}{\varepsilon}

\title{Concurrency Theory --- Tutorial 1\\[2mm]
\large Transition Systems and Operational Models}
\author{}
\date{}

\begin{document}
\maketitle

\noindent\textbf{Goal.} This tutorial introduces a common way to describe the
step-by-step behaviour of a system. We will turn operational rules into a
labelled transition system (LTS), explore the reachable part of that system,
and compare models with different kinds of configurations.

\section*{A small amount of theory}

\paragraph{Labelled transition systems.}
A labelled transition system is a tuple
\[
  \mathcal{T}=(C,\Act,\longrightarrow,c_0),
\]
where $C$ is a set of states (also called configurations), $\Act$ is a set
of action labels, $\longrightarrow\subseteq C\times\Act\times C$ is the
transition relation, and $c_0\in C$ is the initial state. We write
\[
  c\xrightarrow{a}c'
\]
when the system can move from $c$ to $c'$ by performing action $a$.
The relation need not be deterministic: several transitions with the same
source and label may be possible.

A finite run is a sequence
\[
  c_0\xrightarrow{a_0}c_1\xrightarrow{a_1}\cdots
  \xrightarrow{a_{n-1}}c_n.
\]
Its trace is the finite word $a_0a_1\cdots a_{n-1}$. A state is reachable
if some finite run leads to it. A deadlock is a state with no outgoing
transitions.

\paragraph{Operational models and induced LTSs.}
An operational model describes what a system stores and which one-step
updates are allowed. Its configurations contain all information needed to
determine the available next steps. The \emph{induced LTS} has:
\begin{itemize}
  \item one LTS state for each configuration;
  \item an edge $c\xrightarrow{a}c'$ for each allowed one-step update
        from $c$ to $c'$ labelled by $a$;
  \item the model's initial configuration as the initial LTS state.
\end{itemize}
The reachable part of the induced LTS contains just the configurations
reachable from the initial one. It can be finite or infinite.

\paragraph{A useful graph convention.}
When discussing treewidth, we forget edge directions and labels: two
configurations are adjacent in the underlying undirected graph if there is
an LTS transition between them in either direction.

\paragraph{Forward and backward reachability.}
For a transition relation $\longrightarrow$ and a set of configurations
$X$, define
\[
      \operatorname{Post}^*(X)=\{c': c\in X\text{ and }c\longrightarrow^*c'\},
      \qquad
      \operatorname{Pre}^*(X)=\{c: c\longrightarrow^*c'\text{ for some }c'\in X\}.
\]
Here $\longrightarrow^*$ allows zero or more steps. Thus
$\operatorname{Post}^*(\{c_0\})$ is the reachability set from the initial
configuration, while $\operatorname{Pre}^*(F)$ contains configurations
from which some target in $F$ can be reached.

\paragraph{Semilinear sets.}
A linear set in $\mathbb{N}^d$ has the form
\[
      \{b+n_1p_1+\cdots+n_kp_k:n_1,\ldots,n_k\in\mathbb{N}\},
\]
for fixed vectors $b,p_1,\ldots,p_k\in\mathbb{N}^d$. A semilinear set is
a finite union of linear sets. For one counter, semilinear subsets of
$\mathbb{N}$ are exactly the ultimately periodic sets; equivalently, when
counter values $n$ are encoded by $1^n$, they are exactly the regular unary
languages. For stacks, the analogous notion we will use is regularity of
the stack words: a set of pushdown configurations is regular if, for each
control state, its stack contents form a regular language.

The right representation depends on the model. Counter values are numbers,
whereas pushdown configurations contain words. For a context-free grammar,
configurations are sentential forms; one can ask whether their language is
regular or context-free, and separately whether its Parikh image is
semilinear.

\section*{Exercises}

\subsection*{1. Reading a finite LTS}
Consider the LTS with states $C=\{r,g\}$, initial state $r$, and transitions
\[
  r\xrightarrow{\mathsf{go}}g,
  \qquad
  g\xrightarrow{\mathsf{stop}}r,
  \qquad
  g\xrightarrow{\mathsf{wait}}g.
\]
\begin{enumerate}
  \item Draw the LTS and mark its initial state.
  \item List the reachable states, all deadlocks, and the traces of length
        at most three.
  \item Add one transition so that the system has a deadlock. Give two
        different choices and explain the difference between them.
\end{enumerate}

\subsection*{2. A one-counter system}
A one-counter system has a finite set of control states $Q$ and a counter
in $\mathbb{N}=\{0,1,2,\ldots\}$. A rule
\[
  p\xrightarrow{a,d}q
\]
changes the control state from $p$ to $q$, performs action $a$, and changes
the counter by $d\in\mathbb{Z}$. From configuration $(p,n)$ it induces
\[
  (p,n)\xrightarrow{a}(q,n+d)
\]
exactly when $n+d\geq 0$.

Take $Q=\{q\}$, initial configuration $(q,0)$, and rules
\[
  q\xrightarrow{\mathsf{inc},+1}q,
  \qquad
  q\xrightarrow{\mathsf{dec},-1}q.
\]
\begin{enumerate}
  \item Draw the reachable part of the induced LTS up to counter value $3$.
        Make clear which edges continue beyond the part shown.
  \item Is $(q,0)$ a deadlock? Which actions are enabled at $(q,n)$?
  \item Now allow the counter to range over $\mathbb{Z}$ and remove the
        non-negativity condition. Draw the part with counter values from
        $-2$ to $2$. Compare the two systems: give a finite trace possible
        over $\mathbb{Z}$ but impossible over $\mathbb{N}$.
  \item Is the reachable LTS finite in either version? Explain briefly.
\end{enumerate}

\subsection*{3. A pushdown system}
A pushdown system has a finite set of control states and a stack over a
finite alphabet. The rightmost symbol is the top of the stack. For this
exercise, let the stack alphabet be $\{\bot,0,1\}$, where $\bot$ is a
permanent bottom symbol. There is one control state $q$, the initial stack
is $\bot$, and the rules are:
\begin{itemize}
  \item push a $0$ or a $1$ above the current top, with action
        $\mathsf{push}_0$ or $\mathsf{push}_1$;
  \item pop the top symbol with action $\mathsf{pop}$ whenever the top is
        $0$ or $1$; $\bot$ is never popped.
\end{itemize}
For example, pushing $1$ from stack $\bot 0$ gives
$(q,\bot 0)\xrightarrow{\mathsf{push}_1}(q,\bot 01)$.
\begin{enumerate}
  \item Draw the reachable induced LTS through stack height three,
        labelling push and pop edges. (The height is the number of symbols
        above $\bot$.)
  \item Describe the underlying undirected graph. What familiar tree does
        it resemble, and how does the root's degree differ from that of
        other vertices?
  \item Add a second control state $r$ and an action $\mathsf{switch}$ that
        changes $q$ to $r$ and back without changing the stack. What changes
        in the first two levels of the induced LTS?
  \item Which information must be included in a configuration? Why would
        recording only the stack, or only the control state, lose
        information about future behaviour?
\end{enumerate}

\subsection*{4. Independent steps and interleaving}
Consider a system with configuration $(x,y)\in\{0,1\}^2$, initially
$(0,0)$. Action $a$ changes $x$ from $0$ to $1$ and leaves $y$ unchanged;
action $b$ changes $y$ from $0$ to $1$ and leaves $x$ unchanged. Each
action can be performed only once.
\begin{enumerate}
  \item Construct the induced LTS. Label every state by its pair $(x,y)$.
  \item Show that both traces $ab$ and $ba$ are possible and lead to the
        same configuration. Draw the resulting diamond.
  \item Identify the two enabled actions at the initial configuration.
        What does this example suggest about representing concurrency by
        an LTS whose runs execute one action at a time?
  \item Change the rules so that $b$ is enabled only after $a$. Which of
        the two traces remains possible?
\end{enumerate}

\subsection*{5. From rules to an induced LTS}
For each model in Exercises 2--4, explicitly identify its configuration
space, initial configuration, action labels, and one-step transition
relation. Then compare the resulting reachable LTSs:
\begin{enumerate}
  \item Which have finitely many reachable configurations?
  \item Which have a branching structure, and which have cycles?
  \item Give an example of two different configurations with the same
        enabled action labels. Does that imply that their future behaviours
        are identical? Justify your answer with a path or counterexample.
\end{enumerate}

\subsection*{6. Reachability sets: semilinear or regular?}
In this exercise, a one-counter net has finite control and counter updates
by fixed integers, enabled when the resulting counter remains
nonnegative. A one-counter automaton additionally has guards testing
whether the counter is zero or positive. For both models, identify a
configuration $(q,n)$ with its control state and counter value.
\begin{enumerate}
  \item For a one-counter net, prove or refute that both
        $\operatorname{Post}^*(X)$ and $\operatorname{Pre}^*(X)$ are
        semilinear whenever $X$ is semilinear.
  \item Repeat part 1 for one-counter automata with zero tests. State the
        corresponding regular-language formulation using the unary encoding
        $1^n$. In both parts, note that a finite initial set is a special
        case of a semilinear set.
  \item For a pushdown system, take a regular set $X$ of configurations.
        Show that both $\operatorname{Post}^*(X)$ and
        $\operatorname{Pre}^*(X)$ are regular sets of configurations. What
        does this imply for a DPDA, viewed as a pushdown system? Is
        determinism needed for the reachability-set result?
  \item Let $G$ be a context-free grammar with start variable $S$ and
        terminal alphabet $T$. Encode a sentential form over $T\cup V$ by
        treating each variable as a distinct output symbol. Prove that
        $\operatorname{Post}^*(\{S\})$ is context-free: for each variable
        $A$, add a production $A\to\widehat{A}$, where $\widehat{A}$ is a
        fresh terminal symbol meaning ``leave this variable unexpanded.''
        Deduce that the Parikh image is semilinear. Must this language of
        sentential forms be regular?
  \item For the same grammar, let $K\subseteq T^*$ be a regular language
        of terminal target words. Show that the sentential forms in
        $\operatorname{Pre}^*(K)$ form a regular language over $T\cup V$.
        Hint: take a DFA for $K$ and, for each grammar variable $A$, record
        the pairs of DFA states connected by some word generated from $A$.
\end{enumerate}

\subsection*{7. Same language, different transition systems: BPA and PDA}
Language equivalence compares only which finite words are accepted; it
does not require the induced LTSs to have the same structure. Consider
the BPA process given by the recursive equation
\[
  X = 1 + a.X + b.X,
\]
where $1$ denotes successful termination. Its induced LTS has one
configuration $X$ with transitions $X\xrightarrow{a}X$ and
$X\xrightarrow{b}X$. Its language consists of the words labelling runs
that end successfully.

Now consider the PDA with initial control state $q$, permanent stack
symbol $\bot$, and accepting states $q$ and $r$. It never changes the
stack. From $q$ it has transitions labelled $a$ and $b$ back to $q$, and
an additional transition labelled $a$ from $q$ to $r$. State $r$ has no
outgoing transitions. A PDA accepts a word if some run consumes the whole
word and ends in an accepting state.
\begin{enumerate}
  \item Write a context-free grammar for the language of the BPA process
        and determine the language accepted by the PDA. Show that both
        languages are $\{a,b\}^*$.
  \item Draw the two induced LTSs, including the successful-termination
        or accepting-state information. Are their transition graphs
        isomorphic?
  \item Prove that the initial states are not strongly bisimilar. In
        particular, consider the PDA's $a$-transition from $q$ to $r$:
        which BPA state could match it, and why can the match not be
        continued?
  \item Explain why this example does not contradict the equivalence
        between context-free grammars and PDAs: what information does
        language equivalence forget?
\end{enumerate}

\subsection*{8. Challenge: a nonsemilinear Turing-machine reachability set}
For this exercise, use the following explicit convention. A Turing-machine
configuration has a finite-control state and a tape encoding. At selected
checkpoint configurations, the tape has the form
\[
  \vdash\, a^n\#b^m\,\dashv,
\]
with the head at a fixed location. The checkpoint projection records
$(n,m)$. A set of configurations is called semilinear here if its
projection to the control state and these symbol-count coordinates is
semilinear; in particular, every fixed-control-state slice and every
coordinate projection of such a set must be semilinear.
\begin{enumerate}
  \item Design a deterministic Turing machine that reaches checkpoint
        configurations with tape contents
        $\vdash\,a^n\#b^{n^2}\,\dashv$ for every $n\geq 1$, and enters
        the checkpoint state only for these canonical configurations.
        Sketch how it can maintain the unary value $n$ and construct the
        block of $n^2$ symbols.
  \item Prove that $\{(n,n^2):n\geq 1\}$ is not semilinear. Use the
        checkpoint slice to conclude that the projected reachability set
        of your machine is not semilinear.
  \item Why is the encoding/projection convention important? Explain why
        ``semilinear'' is not a representation-free property of arbitrary
        sets of Turing-machine configurations.
\end{enumerate}

\subsection*{9. Optional challenge: treewidth of a pushdown graph}
This challenge uses the following normal form. Let $Q$ be a finite set of
control states and $\Gamma$ a finite stack alphabet containing $\bot$.
Every transition replaces the top stack symbol $A$ by a word $\beta$ of
length at most two, so that
\[
  (p,wA)\longrightarrow(q,w\beta),
  \qquad |\beta|\leq 2,
\]
and no transition removes $\bot$. Consider the underlying undirected graph
of all configurations reachable from a fixed initial configuration.

A tree decomposition consists of a tree of bags of vertices such that
every vertex and edge occurs in some bag, and the bags containing any
fixed vertex form a connected subtree. The width is the largest bag size
minus one.
\begin{enumerate}
  \item Use the tree of stack words as a decomposition tree. For each word
        $w\in\Gamma^*$, consider the bag
        \[
          B_w=(Q\times\{w\})\;\cup\;
              \bigcup_{A\in\Gamma}(Q\times\{wA\}).
        \]
            Intersect each bag with the vertex set of $G$ if $G$ contains only
            reachable configurations.
        Explain why every configuration occurs in at least one bag and
        why the bags containing a fixed configuration are connected.
  \item Check that every possible transition edge has both endpoints in
        a common bag. Separate the cases where the stack height stays the
        same, decreases by one, or increases by one.
  \item Deduce the bound
        \[
          \operatorname{tw}(G)\leq |Q|(|\Gamma|+1)-1.
        \]
\end{enumerate}

\paragraph{Further questions for discussion.}
For a first meeting, useful habits beyond drawing the graph are to state
the configuration precisely, distinguish the full LTS from its reachable
part, check when actions are enabled, find deadlocks and cycles, compare
traces, and look for simple invariants. When two actions appear independent,
check whether both execution orders are possible and reach the same state.
These checks apply to finite-state systems, counters, stacks, and token
systems alike.

\end{document}